\chapter{Methodology}
\label{ch:methodology}

Our simulation framework models a stylized market environment in which agents
repeatedly allocate resources.  This approach allows us to computationally
test theoretical trade-offs under rigorous, repeatable conditions.

% ---------------------------------------------------------------
\section{The Stylized Market Environment}
% ---------------------------------------------------------------

We model the market as a filtered probability space
$(\Omega, \Filt, \{\Filt_t\}, \Prob)$ (Definition~2.1) in which returns
$r_t \mid S_t$ are drawn from one of $M$ regime-specific distributions.

\subsection{Regime Configurations}

Three synthetic environments are tested:

\paragraph{Stationary.}  A single regime with
$r_t \sim \Normal(\mu, \sigma^2)$, $\mu = 0.08$, $\sigma = 0.15$.
This environment tests pure exploitation and validates the SLLN result of
Theorem~\ref{thm:kelly_slln}.

\paragraph{Adversarial Shock.}  A deterministic structural break at step
$T_{\text{shock}} = T/2$: Bull market
$r_t \sim \Normal(0.08, 0.15^2)$ for $t \leq T_{\text{shock}}$, then
Bear market $r_t \sim t_3(-0.10, 0.30^2)$ (Student-$t$ with three degrees
of freedom) for $t > T_{\text{shock}}$.  This tests robustness to sudden
regime change and replicates the Kelly-Ruin Paradox.

\paragraph{Slow Gaussian Drift.}  The mean drifts linearly:
$\mu_t = \mu_0 + (\mu_1 - \mu_0) \cdot t / T$, with $\mu_0 = 0.08$ and
$\mu_1 = -0.05$.  This tests the ability of agents to unlearn their prior
gradually.

\subsection{Common Random Numbers (CRN)}

All agents in a given experiment are evaluated on the \emph{same} realized
return matrix $\mathbf{R} \in \R^{N_{\text{sim}} \times T}$, generated once
before any agent acts.  This is the Common Random Numbers (CRN) variance
reduction technique (L'Ecuyer, 1994): by controlling for Monte Carlo noise,
differences in terminal wealth between agents are attributable solely to
strategy differences, not sampling variation.  Any comparison that violates
CRN is statistically invalid.

% ---------------------------------------------------------------
\section{Decision Agents}
% ---------------------------------------------------------------

\subsection{Kelly Oracle (Benchmark)}

The Oracle knows the true distribution parameters at every step and computes
$f^* = \mu/\sigma^2$ (Equation~\eqref{eq:gaussian_kelly}).  It represents the
theoretical upper bound on performance and is used to compute the \emph{oracle
regret} $\mathcal{R}_t^{\text{oracle}} = \log W_t^{(f^*)} - \log W_t^{(\text{agent})}$
plotted in the regret dynamics charts.

\subsection{Naive Bayes Kelly}

Maintains Beta posteriors $(\alpha_k, \beta_k)$ and plugs the posterior mean
$\hat p_k = \alpha_k / (\alpha_k + \beta_k)$ directly into the Kelly formula.
As formalized in Proposition~\ref{prop:sticky_prior}, this agent suffers
$O(T_0)$ detection lag after a regime change because its posterior
concentrates mass that resists reallocation.

\subsection{Thompson Sampling Kelly}

Samples $\tilde p_k \sim \text{Beta}(\alpha_k, \beta_k)$ and bets
$f^*(\tilde p_k)$.  The Bayesian posterior variance acts as an implicit
risk aversion: during the first $O(\sqrt{T})$ steps when variance is high,
the sampled $\tilde p_k$ are dispersed, leading to lower average bets and
thus lower early-stage ruin risk compared to the Naive plug-in.

\subsection{UCB Agent}

Uses the UCB1 rule (Theorem~\ref{thm:ucb_regret}) with the Kelly formula
applied to the upper-confidence-bound estimate of each arm's win probability.

\subsection{EXP3 Agent (Fractional Variant)}

\footnote{See Remark~\ref{rem:exp3_variant} in Chapter~\ref{ch:foundations}.
Our implementation is a continuous-fraction variant of EXP3 and does not
satisfy the adversarial regret bound of Theorem~\ref{thm:exp3_regret}
exactly; we use it as an adversarially-motivated heuristic rather than
claiming the exact theoretical guarantee.}
Maintains exponential weights $w_k$ updated via
$w_k \leftarrow w_k \exp(\gamma \hat r_k / \Prob(k))$
and allocates fractions proportional to the resulting mixing distribution.

\subsection{Proposed Method 1: Volatility-Augmented HMM (Vol-HMM)}

The Vol-HMM addresses the Sticky Prior Paradox through three mechanisms:

\paragraph{Two-dimensional observation.}
At each step, the agent observes $(r_t, v_t)$ where
$v_t = \hat\sigma_{[t-W, t]}$ is the rolling standard deviation over a window
$W = 5$.  Because $v_t$ spikes immediately when the market shifts from low-
to high-volatility regimes, it provides an early-warning signal that the
return alone cannot.

\paragraph{Log-space forward algorithm.}
Regime beliefs are maintained via Equation~\eqref{eq:forward_logspace}.
The joint emission log-likelihood is
\begin{equation}
    \ln b_j(r_t, v_t) \;=\;
    \ln \mathcal{N}(r_t;\,\mu_j,\,\sigma_j^2)
    + \ln \Gam(v_t;\,k_j,\,\theta_j),
    \label{eq:vol_hmm_emission}
\end{equation}
where $\mathcal{N}$ is the Gaussian density and $\Gam$ is the Gamma density.
Emission parameters are specified a priori (Remark~\ref{rem:hmm_specified})
as $(\mu_{\text{bull}}, \sigma_{\text{bull}}, k_{\text{bull}},
\theta_{\text{bull}}) = (0.08, 0.15, 2.0, 0.075)$ and
$(\mu_{\text{bear}}, \sigma_{\text{bear}}, k_{\text{bear}},
\theta_{\text{bear}}) = (-0.10, 0.30, 4.0, 0.15)$.

\paragraph{CUSUM change-point trigger.}
In addition to the Bayes filter, a CUSUM statistic accumulates evidence of
a negative mean shift:
\begin{equation}
    S_t^- \;=\; \max\!\bigl(0,\; S_{t-1}^- - z_t - k_{\text{drift}}\bigr),
    \qquad z_t = (r_t - \mu_{\text{bull}}) / \sigma_{\text{bull}}.
    \label{eq:cusum}
\end{equation}
When $S_t^- > h$ (threshold $h = 3.0$), a +2 log-likelihood bonus is added
to the bear state emission, accelerating transition of belief mass.

\paragraph{Transition constraint.}
The transition matrix satisfies $A_{ii} \leq 0.95$ to prevent the sticky-
prior effect from re-emerging through the transition dynamics.

\paragraph{Bet sizing.}
The bet fraction is $f_t = P(\text{bull} \mid \Filt_t) \cdot f^*_{\text{bull}}$,
reducing exposure linearly as probability mass flows to the bear state.

\subsection{Proposed Method 2: Risk-Constrained Kelly (CPPI)}

\begin{definition}[CPPI Floor and Cushion]
Let $\bar W_t = \max_{s \leq t} W_s$ be the running peak wealth.  The
\emph{floor} is $F_t = (1 - D_{\max})\bar W_t$ and the \emph{cushion} is
$C_t = W_t - F_t$.
\end{definition}

\begin{proposition}[CPPI Leverage Bound]
\label{prop:cppi_bound}
Under the leverage rule $\lambda_t = \min(1, m \cdot C_t / W_t)$ with
$m > 0$, the wealth process satisfies
\[
    W_{t+1} \;\geq\; F_t \cdot \bigl(1 - m\,C_t\,r_{\min} / W_t\bigr),
\]
where $r_{\min}$ is the worst-case single-period return.  In continuous
time (or discrete time with $|r_t| \leq m^{-1}$), this guarantees
$W_t \geq F_t$ for all $t$, i.e., the drawdown never exceeds $D_{\max}$.
\end{proposition}

\begin{proof}
The wealth update is
$W_{t+1} = W_t(1 + \lambda_t r_t)$.
With $\lambda_t = m C_t / W_t$ (ignoring the $\min$ cap),
$W_{t+1} = W_t + m C_t r_t$.
In the worst case $r_t = r_{\min} < 0$:
$W_{t+1} = W_t + m C_t r_{\min} \geq F_t + C_t(1 + m r_{\min}) \geq F_t$
provided $m |r_{\min}| \leq 1$.  The discrete-time caveat is that a single
return $r_t < -1/m$ can breach the floor; this is noted in the experimental
discussion.
\end{proof}

The CPPI agent does not use Bayesian updating of return parameters; it is a
purely mechanical risk-control layer applicable to any underlying Kelly
strategy.  We use the posterior mean as the base Kelly fraction and apply
the CPPI leverage multiplier on top.

% ---------------------------------------------------------------
\section{Market Friction and Transaction Costs}
% ---------------------------------------------------------------

To bridge the gap between theoretical growth and empirical deployment, we
model transaction costs at each step.  Given betting fraction $f_t$ and the
\emph{drift-adjusted} previous fraction
$\hat f_{t-1} = f_{t-1}(1+r_{t-1}) / (1 + f_{t-1} r_{t-1})$, the
cost is applied to turnover:
\begin{equation}
    W_{t+1} \;=\; W_t \cdot \bigl(1 + f_t r_t - c\,|f_t - \hat f_{t-1}|\bigr),
    \label{eq:friction}
\end{equation}
where $c$ is the friction in basis points ($\times 10^{-4}$).  We test
$c \in \{0, 10, 25, 50, 100\}$ bps.  High-turnover agents (EXP3, UCB)
rebalance aggressively and thus incur proportionally higher fees.

% ---------------------------------------------------------------
\section{Evaluation Metrics}
% ---------------------------------------------------------------

Strategies are evaluated against four complementary objectives:

\begin{itemize}
    \item \textbf{Median Terminal Wealth} $\tilde W_T$: The 50th percentile
    of the terminal wealth distribution across $N_{\text{sim}}$ paths.  We
    use the median rather than the mean because the distribution is
    log-normal with a heavy right tail.

    \item \textbf{Probability of Ruin}:
    $\hat\rho = N_{\text{ruined}} / N_{\text{sim}}$, estimated via Monte
    Carlo.  Statistical significance is assessed by Fisher's Exact Test
    on the $2 \times 2$ contingency table of ruin/survival counts.

    \item \textbf{Annualized Sharpe Ratio}:
    $\mathcal{S} = (\bar R / \hat\sigma_R) \cdot \sqrt{252}$,
    computed via the $O(T)$ running-statistics recursion
    $\bar R_t = \bar R_{t-1} + (R_t - \bar R_{t-1})/t$.

    \item \textbf{Annualized Sortino Ratio}:
    $\mathcal{SOR} = (\bar R / \hat\sigma_{\downarrow}) \cdot \sqrt{252}$,
    where $\hat\sigma_{\downarrow}^2 = \frac{1}{T}\sum_t \min(R_t, 0)^2$.
\end{itemize}

The cumulative \emph{Oracle Regret} $\log W_t^{(f^*)} - \log W_t^{(\text{agent})}$
is plotted across time to reveal the dynamic rate at which each agent
converges to or diverges from the growth-optimal benchmark.
